\documentclass[11pt]{article}
\usepackage[T1]{fontenc}
\usepackage[utf8]{inputenc}
\usepackage[margin=1.1in]{geometry}
\usepackage{amsmath,amssymb,amsthm}
\usepackage{booktabs,alltt}
\usepackage{hyperref}

\newcommand{\lean}[1]{\texttt{\detokenize{#1}}}
\newcommand{\repo}{\url{https://github.com/unbalancedparentheses/lean-extremevaluetheory}}
\newcommand{\release}{\texttt{v0.9.0}}
\newcommand{\doi}{\texttt{[DOI: fill in after the release]}}

\title{A library of extreme value theory in Lean 4}
\author{Federico Carrone}
\date{November 2026}

\begin{document}
\maketitle

\begin{abstract}
We describe a Lean~4 library, built on Mathlib, that formalizes the core of extreme value theory:
the Fisher--Tippett--Gnedenko theorem, domains of attraction, regular variation, the
Pickands--Balkema--de Haan theorem, estimators and likelihood theory for the tail, stable laws and
heavy tails, extremes of dependent sequences, multivariate extremes and point processes. The library
has about 100,000 lines and 4,300 theorems and lemmas, uses no \texttt{sorry} and depends only on Lean's three standard
axioms, which continuous integration checks for every declaration. Since a formal proof does not
show that a statement matches the literature, the library also carries specification tests, mutation
checks and written audits, and it records where its statements are narrower than their textbook forms.
Formalizing the extremal index exposed its hypotheses and led to examples that, to our knowledge, are
new, including a process that answers a question of Bai, Kulik and Wang. The statements, tests and
audits were written with AI assistance and have not yet been reviewed by human experts in extreme value theory.
\end{abstract}

\section{Introduction}

Extreme value theory describes the behaviour of maxima and of the far tail of a distribution. It is
used to estimate quantities such as the level of a flood expected once in a hundred years or the loss
from a catastrophic insurance event. Its results carry many technical hypotheses, and
software such as \texttt{evd}, \texttt{ismev}, \texttt{extRemes} and \texttt{pyextremes} relies on them.
A machine-checked account makes those hypotheses explicit.

This paper describes such an account, in Lean~4~\cite{lean4} on top of Mathlib~\cite{mathlib}. The
development of the Fisher--Tippett--Gnedenko theorem and its blueprint are based on
ExtremeValueProject, by Kalle Kytölä and the participants of an Aalto University course on formalized
mathematics~\cite{evp}. The source, an interactive blueprint, API documentation and a catalogue of
results with their references are at \repo; this paper refers to release \release{} (\doi).

Our main points are the following.
\begin{itemize}
  \item The library's scope (Section~\ref{sec:contents}), which covers the standard textbooks'
        theory of maxima, regular variation and peaks over threshold, and much of their statistics.
  \item How the library tests that a statement is the right one, and not only that its proof is
        correct (Section~\ref{sec:trust}).
  \item What the process found (Section~\ref{sec:new}): hypotheses in the literature that cannot be
        dropped, and new examples for the extremal index.
  \item What it does not establish (Section~\ref{sec:limits}).
\end{itemize}

\section{What is formalized}\label{sec:contents}

The central result is the classification of limits of normalized maxima. Writing
$A\bullet G$ for the distribution function $G$ after an increasing affine change of variables, the
library states it as follows.
\begin{alltt}\small
theorem CumulativeDistributionFunction.IsExtremeValueDistr.classification
    \{G : CumulativeDistributionFunction\} (hG : G.IsExtremeValueDistr) :
    (\(\exists\) (A : AffineIncrEquiv), A \(\bullet\) G = standardGumbelCDF) \(\lor\)
      (\(\exists\) (A : AffineIncrEquiv) (\(\xi\) : \(\mathbb{R}\)) (h\(\xi\) : 0 < \(\xi\)),
        A \(\bullet\) G = standardFrechetCDF h\(\xi\)) \(\lor\)
      (\(\exists\) (A : AffineIncrEquiv) (\(\xi\) : \(\mathbb{R}\)) (h\(\xi\) : 0 < \(\xi\)),
        A \(\bullet\) G = standardWeibullCDF h\(\xi\))
\end{alltt}
Here \lean{G.IsExtremeValueDistr} says that $G$ is non-degenerate and is a limit of $F(a_n x + b_n)^n$
at continuity points of $G$, for some distribution function $F$ and affine normalizations.
Table~\ref{tab:contents} summarizes the rest; the repository's \texttt{THEOREMS.md} lists each result with its source.

\begin{table}[ht]
\centering\small
\begin{tabular}{p{0.24\linewidth}p{0.68\linewidth}}
\toprule
Area & Main results \\
\midrule
Limits of maxima & Fisher--Tippett--Gnedenko; max-stable laws are exactly the extreme value laws \\
Domains of attraction & Fr\'echet, Weibull and Gumbel characterizations; the quantile form \\
Peaks over threshold & Pickands--Balkema--de Haan as an equivalence for every shape \\
Regular variation & Karamata (both halves), Potter bounds, representation theorem, Breiman's lemma \\
Stable laws & Generalized CLT for $0<\alpha<2$; domains of attraction; uniqueness of the index \\
Heavy tails & Subexponential laws, random sums, Embrechts--Veraverbeke ruin asymptotics \\
Estimators & Hill (consistency, normality, dependent data), Pickands and moment consistency, Weissman \\
Likelihood & GPD and GEV maximum likelihood: consistency, Fisher information; GPD joint normality \\
Return levels & Asymptotic normality, Wald intervals with asymptotic coverage, the delta method \\
Dependence & Leadbetter's theorem; extremal index; tail process; $\theta=\vartheta$ under anticlustering and $A(u_n)$ \\
Time series & AR(1), ARMAX, linear processes, Markov chains; $X=AX+B$ and GARCH(1,1) tail indices \\
Multivariate & Regular variation, converse Cram\'er--Wold, spectral representation, the multivariate GPD \\
Processes & Poisson random measures; exceedance point processes; $D[0,1]$ is Polish; tail empirical process \\
\bottomrule
\end{tabular}
\caption{Contents of the library. Section~\ref{sec:limits} qualifies these entries.}
\label{tab:contents}
\end{table}

General results that are not specific to extremes (vague convergence, Poisson random measures, mixing
coefficients, the Skorokhod space, the delta method) were absent from Mathlib when we needed them. They
live in a separate folder, written to Mathlib's conventions so that they can be contributed upstream.

\section{Testing the statements}\label{sec:trust}

Lean's kernel guarantees that a conclusion follows from its hypotheses. It does not guarantee that the
hypotheses can hold, or that a definition means what its name suggests. In a formal development these
failures are easy to produce: in Lean $x/0=0$, $\log 0=0$, and the integral of a non-integrable
function is $0$, so a statement can be true for an unintended reason. Several statements in the
project this library started from were false as written and have been corrected. We use the following
checks, all but the last of which can be rerun from the repository.
\begin{description}
  \item[Axiom audit.] Continuous integration fails if any of the library's 9,487 declarations depends on an
    axiom other than \lean{propext}, \lean{Classical.choice} and \lean{Quot.sound}.
  \item[Independent kernel.] A nightly job rechecks all declarations, including those of
    Mathlib, with \texttt{con-leche}, a second kernel.
  \item[Non-vacuity tests.] 92 tests apply a headline theorem to a concrete object, showing that its
    hypotheses can hold together. Three of the results in Section~\ref{sec:new} do not have such a test yet.
  \item[Known-answer tests.] 159 tests check definitions against textbook values, which catches sign
    and parametrization errors.
  \item[Junk-value audit.] A scan, run by hand, of 3,718 theorems for statements that hold only through a totalization
    convention found one such statement, now fixed.
  \item[Mutation checks.] 323 perturbations of the statements of 45 headline theorems were tried by hand with their proofs kept.
    A perturbation that still compiles suggests a theorem proved for a trivial reason. One survived, and we
    judged it harmless.
  \item[Faithfulness audits.] Six written comparisons of statements with the literature, recording each deviation.
\end{description}
These tests and audits were themselves written with AI assistance, so they are not independent of the
statements they test. They reduce the risk of error, but do not remove it.

\section{What the formalization found}\label{sec:new}

\paragraph{Hypotheses that cannot be dropped.}
The library contains counterexamples showing that hypotheses of Karamata's theorem, Breiman's lemma, the
von Mises condition, the generalized CLT, Hill consistency, Leadbetter's theorem and the
identification $\theta=\vartheta$ of the extremal index are needed.

\paragraph{The extremal index and the tail process.}
For a stationary regularly varying sequence let $\theta$ be its extremal index, and let $\vartheta$ be the
candidate extremal index read off the spectral tail process~\cite{ks20}. Then $\theta\le\vartheta$, and they agree
under anticlustering together with the block condition $A(u_n)$ on a common block sequence. Bai and Wang~\cite{bw} gave a model with $\theta<\vartheta$. We give a much simpler one.

The \emph{echo process} is built from iid Pareto$(\alpha)$ shocks. Each shock of size $z$ appears
once at its own time and again after a lag that grows with $z$. Since the lag grows without bound, no
fixed lag sees the echo, so the spectral tail process is $\Theta_t=\mathbf 1\{t=0\}$ and $\vartheta=1$. Yet large values
still come in pairs, and the extremal index is $\theta=1/2$
(\lean{logEchoProcess_theta_lt_vartheta}).

\paragraph{An answer to a question of Bai, Kulik and Wang.}
Bai, Kulik and Wang~\cite{bkw} ask whether there are regularly varying processes for which anticlustering
holds with $\lim\log r_n/\log n=\rho$ only for $\rho$ up to some $\rho_0<1$. Taking the lag $\lfloor z^\gamma\rfloor+1$
with $0<\gamma<\alpha$ gives such a process for every $\rho_0=\gamma/\alpha\in(0,1)$: anticlustering at the normalizing levels holds if
$\rho<\rho_0$ and fails if $\rho>\rho_0$ (\lean{powEchoProcess_antiClustering_iff}).

\paragraph{Mixing and a warning for estimators.}
The echo processes are $\beta$-mixing, geometrically for the logarithmic lag, and still have $\theta\neq\vartheta$
(\lean{logEchoProcess_betaMixing_incompatible}). This illustrates the classical point that anticlustering and mixing must
hold at a common block length. On the same process, the runs estimator with short runs converges to
$\vartheta$, not to $\theta$ (\lean{logEcho_tendstoInMeasure_runsEstimator}).

\paragraph{Status of these claims.}
Our literature check compared these results with the papers we could read in full. We could not obtain the Kulik--Soulier
book~\cite{ks20}, Weissman--Cohen, or Smith (1988), so the claim of novelty is provisional, and no expert has confirmed it.
Some results close to these are known: $\theta\le\vartheta$ is Lemma 7.5.4 of~\cite{ks20} as cited by~\cite{bkw},
and the principle behind the mixing example is classical.

\section{Limitations}\label{sec:limits}

Many results are narrower than their textbook forms. The README and \texttt{FAITHFULNESS.md} give the full list; the most important are the following.
\begin{itemize}
  \item \textbf{No expert review.} This matters most, and the rest of this list may be incomplete because of it.
  \item \textbf{GEV likelihood.} Maximum likelihood consistency is proved on a compact parameter set and for an exact iid GEV sample, not for block maxima from a law in a domain of attraction.
    Joint normality of the three-parameter estimator is not proved, so GEV and Gumbel return-level normality
    assume parameter normality as a hypothesis.
  \item \textbf{GPD likelihood and return levels.} Joint normality needs $\xi_0>-1/2$, $\xi_0\neq 0$ and a compact parameter
    set. Peaks-over-threshold return levels treat the excesses as an exact iid GPD sample with a known exceedance probability.
  \item \textbf{Estimators.} Normality of the Pickands and moment estimators is not proved, and neither are
    profile-likelihood intervals. The blocks and runs estimators of the extremal index assume that their population ratio converges to $\theta$, and the intervals estimator assumes its empirical-moment limits.
  \item \textbf{Ruin.} The ruin probability is defined by the Pollaczek--Khinchine formula, not derived from the risk process.
  \item \textbf{Stochastic recurrence.} The tail of $X=AX+B$ is identified through moment bounds. The sharp Kesten--Goldie
    asymptotics is not proved.
  \item \textbf{Dependence.} Several theorems assume noise that is nonnegative at every sample point.
  \item \textbf{Asymptotics only.} The library says nothing about finite-sample behaviour or about choosing a threshold or block length.
\end{itemize}

\section{How it was built}

Most of the code, tests and documents were written by AI coding agents (Claude Code), directed and reviewed by the author, on top of the
Aalto project. Codex, a second AI system, reviewed the statements and claims of this release. This changes what evidence the library
can offer. The proofs need no trust in the agents, because the kernel checks them. The statements do, and that is the reason for
Section~\ref{sec:trust} and for the request below.

\section{Conclusion}

We would most value review by extreme value theorists. The repository's \texttt{docs/REVIEW.md} translates each statement into ordinary mathematics,
so no knowledge of Lean is needed, and corrections are very welcome. Next come the missing likelihood results, contributing the general
results to Mathlib, and executable estimators cross-checked against R and Python packages.

\begin{thebibliography}{9}
\bibitem{bkw} S.~Bai, R.~Kulik and Y.~Wang. A remarkable example on clustering of extremes for regularly-varying stochastic processes. arXiv:2409.17966; Ann.\ Appl.\ Probab.\ 36(4) (2026) 3132--3168.
\bibitem{bw} S.~Bai and Y.~Wang. Tail processes for stable-regenerative multiple-stable model. arXiv:2110.07499; Bernoulli 29(4) (2023) 3255--3279.
\bibitem{lean4} L.~de Moura and S.~Ullrich. The Lean 4 theorem prover and programming language. CADE 2021.
\bibitem{evp} K.~Kytölä and contributors. ExtremeValueProject. \url{https://github.com/kkytola/ExtremeValueProject}.
\bibitem{ks20} R.~Kulik and P.~Soulier. \emph{Heavy-Tailed Time Series}. Springer, 2020.
\bibitem{mathlib} The mathlib Community. The Lean mathematical library. CPP 2020.
\end{thebibliography}

\end{document}
